\section{Related Work}
\label{sec:related}

\noindent\textbf{Pose estimation and body reconstruction.}
Methods that recover body configuration from a single image treat the geometry as their final product.
Two-dimensional estimators localize joint keypoints accurately.
Successive designs fuse multi-scale features~\cite{newell2016stacked,xiao2018simple}, associate keypoints across multiple people~\cite{cao2019openpose}, and adopt high-resolution or transformer backbones~\cite{sun2019deep,xu2022vitpose}.
Parametric body models~\cite{loper2015smpl} extend this recovery to three-dimensional pose and shape, and regressors estimate both from one image~\cite{kanazawa2018end,kolotouros2019spin,li2022cliff,goel2023humans}.
SAM 3D Body extends this line to reconstruction of the full body~\cite{yang2026sam3dbody}.
We instead treat reconstruction as annotation, converting its dense outputs into interpretable measurements that ground our language supervision.

\noindent\textbf{Action quality and pose correction.}
Action quality assessment judges how well a movement is performed without stating how the body is configured.
These methods score performance across varied actions~\cite{lei2019survey,parmar2019action} and rehabilitation exercises~\cite{liao2020deep}, and yoga systems go further to report the joint angles behind a pose prediction~\cite{dittakavi2022posetutor}.
Those angles explain a verdict on correctness.
Our supervision expresses the configuration in language, with every numeric statement traceable to an explicit, image-specific measurement.

\noindent\textbf{Vision--language models and VQA.}
Vision--language models link appearance to the words that describe it.
Contrastive pretraining over image--text pairs establishes this link at scale~\cite{radford2021clip}.
Later systems couple frozen image encoders with language models~\cite{alayrac2022flamingo,li2023blip2} and tune the pair to follow visual instructions~\cite{liu2023visual,dai2023instructblip,wang2024qwen2vl}.
The benchmarks behind this progress ask about objects, attributes, and relations whose answers are visible in the image~\cite{goyal2017vqav2,hudson2019gqa}.
A body posture, by contrast, is arbitrarily linked to its cultural or social meaning.
Postures that look nearly identical can carry substantially different meanings, and appearance statistics alone cannot determine which meaning applies.
Recovering the meaning therefore requires the measured configuration of the body.

\noindent\textbf{Human pose datasets.}
To our knowledge, no existing pose dataset pairs an image with measured body geometry and language grounded in that geometry.
Large collections annotate two-dimensional keypoints at scale~\cite{lin2014coco,andriluka2014mpii}.
Motion capture supplies three-dimensional pose under controlled conditions~\cite{ionescu2014human36m}, and later work extends it to in-the-wild video~\cite{vonMarcard2018}.
Closest to our setting, Yoga-82 offers a visually diverse benchmark of 82 asanas~\cite{verma2020yoga82}.
Its supervision stops at hierarchical category labels.
YogaAlign builds on that diversity and supplies the missing pairing of reconstruction-derived measurements and language grounded in them.
